\chapter{Background and Literature Review}\label{ch:ch02}

This chapter establishes the intellectual and engineering context in which APEX was
built. It surveys four bodies of work that the system draws on directly: the history and
methods of automated assessment of programming (\Cref{sec:bg-autoassess}); online judge
systems and the sandboxing technology that lets them run untrusted code safely
(\Cref{sec:bg-judges}); learning-management systems and the quiz and coding tooling built
on top of them (\Cref{sec:bg-lms}); and code-similarity detection, which sits adjacent to
grading as a fairness concern (\Cref{sec:bg-plagiarism}). The chapter closes with a gap
analysis (\Cref{sec:bg-gap}) and a comparison table (\Cref{tab:bg-comparison}) that
positions APEX against representative alternatives across the dimensions that matter for a
single-department classroom deployment.

The review is deliberately selective. Its purpose is not to catalogue every assessment
tool ever published, but to show what the mature systems do well, what they trade away,
and where a small self-hostable portal that unifies three answer formats behind one
grading pipeline earns its place. Throughout, claims about APEX are grounded in the actual
implementation rather than aspiration.

\section{Automated Assessment of Programming}\label{sec:bg-autoassess}

\subsection{A short history}

Automatically grading programs is nearly as old as teaching programming. The earliest
graders date to the 1960s, when batch systems compiled a student's submission and compared
its printed output to a reference. The idea has proven remarkably durable: the fundamental
loop---feed an input, capture the output, compare it to what was expected---is still the
core of every modern coding-assessment platform, APEX included.

The first systematic academic survey of the modern era is Douce, Livingstone and
Orwell~\cite{douce2005autoassess}, who in 2005 organised the field into three
``generations'': early command-line graders tied to a single machine; web-based systems
that decoupled the student from the grading host; and a nascent third generation
concerned with pedagogy, feedback quality, and integration with course-management tooling.
Their central observation has aged well: the hard problems in automated assessment are not
the mechanics of running code but the surrounding concerns---security of execution, quality
of feedback, and the integration of grading into an actual course.

Five years later, Ihantola et al.~\cite{ihantola2010autograding} produced the reference
survey for the field, reviewing dozens of tools against a common set of features: supported
languages, the style of test specification, the resubmission policy, the security model,
and---critically---how tightly the grader integrated with a learning-management system.
Their finding was that classroom features and grading features consistently lived in
different products, and that stitching them together fell to the individual instructor.
That observation is precisely the gap APEX targets.

The most cited recent survey, Wa\'{s}ik et al.~\cite{wasik2018survey}, examines online
judges specifically and decomposes them into three recurring components that recur
throughout this chapter: (1) the \emph{contest or exam scaffold} that presents problems and
collects submissions; (2) the \emph{language-aware compile-and-run sandbox} that executes
untrusted code under resource limits; and (3) the \emph{verdict aggregator} that turns
per-test outcomes into a score. APEX maps cleanly onto this decomposition: its scaffold is
the Go/Gin backend and React SPA, its sandbox is delegated to Judge0, and its aggregator is
the \texttt{judge0.Grade} / \texttt{maybeFinalizeAttempt} pipeline analysed in
\Cref{ch:ch07}.

\subsection{Two families of grader: I/O diffing versus unit testing}

The literature divides auto-grading into two broad families, and the choice between them is
the single most consequential design decision a grader makes.

\paragraph{Input/output diffing.} The student writes a complete program that reads from
standard input and writes to standard output. The grader supplies an input, captures the
program's output, and compares it---textually---against an expected output. This is the
model of the competitive-programming judges and the one Douce and Wa\'{s}ik describe as
dominant~\cite{douce2005autoassess,wasik2018survey}. Its virtues are that it is entirely
language-agnostic (the grader never needs to understand the student's code, only its
output), simple to specify (a test case is just a pair of strings), and robust against
students who structure their solution differently than the instructor imagined.

\paragraph{Unit-test execution.} The student writes a function or class conforming to a
fixed signature, and the grader compiles it against an instructor-authored test harness
that asserts on return values. This model, used by tools such as
CodeRunner~\cite{coderunner} and much of the CS-education tooling surveyed by
Ihantola~\cite{ihantola2010autograding}, gives finer-grained feedback (``your function
returned the wrong value for a negative argument'') and can grade code that never produces
console output. Its cost is that the harness is language-specific and the student is
constrained to an exact interface.

APEX uses I/O diffing, for the reasons the survey literature endorses: it keeps the grader
uniform across all six supported languages (Python, JavaScript, TypeScript, Java, C, and
C++) and lets a teacher author a coding problem as nothing more than a list of
\texttt{(input, expected\_output)} pairs. The well-known weakness of I/O diffing is its
brittleness: a solution that computes the correct answer can still be rejected on a
cosmetic mismatch---a trailing newline, a Windows carriage return, or a stray trailing
space. This ``trailing-newline unfairness'' is the most common false negative on any
diff-based platform.

APEX's response is a nine-line normalisation function, \texttt{normalizeOutput}, applied
symmetrically to both the student's actual output and the teacher's expected output before
comparison. It converts \texttt{CRLF} and lone \texttt{CR} line endings to \texttt{LF},
trims trailing spaces and tabs on each line, and trims leading and trailing whitespace of
the whole string. \Cref{lst:bg-normalize} reproduces it. The design intent, developed fully
in \Cref{ch:ch07}, is to make the comparison measure the \emph{correctness of content}
rather than the \emph{byte-exactness of whitespace}. It is honest about its scope: it does
not collapse internal whitespace runs and applies no floating-point tolerance, so it
addresses the most common false negative rather than every conceivable one.

\begin{lstlisting}[language=Go,caption={The output-normalisation function that mitigates I/O-diff brittleness, applied symmetrically to actual and expected output.},label={lst:bg-normalize}]
// normalizeOutput strips CR, trailing whitespace on each line, and trims
// edges so Judge0 output (often \r\n with a trailing newline) compares
// cleanly to teacher-entered expected output.
func normalizeOutput(s string) string {
    s = strings.ReplaceAll(s, "\r\n", "\n")
    s = strings.ReplaceAll(s, "\r", "\n")
    lines := strings.Split(s, "\n")
    for i, ln := range lines {
        lines[i] = strings.TrimRight(ln, " \t")
    }
    return strings.TrimSpace(strings.Join(lines, "\n"))
}
\end{lstlisting}

\subsection{Feedback, partial credit, and resubmission}

Beyond the pass/fail verdict, the assessment literature emphasises three pedagogical
levers. The first is \emph{feedback}: Ihantola et al.\ note that the difference between a
grader that merely scores and one that teaches lies in what it tells the student about a
failure~\cite{ihantola2010autograding}. APEX distinguishes visible \emph{sample} test cases
(\texttt{is\_sample = true}), whose input, expected, and actual output are shown, from
hidden tests, for which the student sees only pass/fail and timing---preventing
reverse-engineering of the graded tests while still offering pre-submission feedback.

The second lever is \emph{partial credit}. Rewarding a solution that passes seven of ten
tests with seventy percent rather than zero is both fairer and more informative. APEX awards
proportional partial credit on coding problems (\texttt{score = passed / total * 100}) but
uses all-or-nothing set-equality on multiple-choice questions, reflecting the different
semantics of the two formats.

The third lever is the \emph{resubmission policy}---how freely a student may iterate.
APEX separates ``run'' from ``submit'': a student may execute sample tests any number of
times during an exam at zero database and grading cost, because only the final
\texttt{POST /student/exams/:id/submit} persists a graded submission. This mirrors the
iterate-freely, grade-once posture that the survey literature associates with good learning
outcomes, without the storage and sandbox load of grading every keystroke.

\section{Online Judges and Sandboxing}\label{sec:bg-judges}

\subsection{The online-judge tradition}

Whereas the assessment literature grew out of education research, the online-judge
tradition grew out of competitive programming, where the correctness and, above all, the
\emph{security} of the execution sandbox dominate every other concern. Systems such as the
UVa Online Judge, Codeforces, and HackerRank~\cite{hackerrank} all share the same shape:
accept a source submission, run it against a battery of tests inside an isolation layer,
and return a verdict. Wa\'{s}ik et al.~\cite{wasik2018survey} survey this lineage in detail,
and C\'{a}ceres et al.~\cite{caceres2014online} survey the web architectures behind such
systems, observing that the durable designs separate the sandbox from the scaffold behind a
thin HTTP interface---exactly the boundary APEX draws between its Go backend and Judge0.

The reason execution is the crux is stark: a student submission is arbitrary, untrusted
code. A single line such as \texttt{os.system("rm -rf /")} or a fork bomb, if run on the
application host, compromises the server. A judge must therefore run every submission inside
a locked environment with strict CPU, memory, and filesystem limits and no meaningful access
to the host or network. Writing such a sandbox correctly is one of the hardest and most
security-sensitive tasks in systems programming, which is why the mature ecosystem has
converged on a small number of shared, battle-tested implementations rather than each
platform rolling its own.

\subsection{The \texttt{isolate} sandbox}

The de facto standard sandbox for this purpose is
\texttt{isolate}~\cite{isolate}, originally developed as the grading sandbox for the
International Olympiad in Informatics (IOI). \texttt{isolate} builds its ``locked room''
from Linux kernel primitives: \emph{namespaces} isolate a process's view of the filesystem,
process table, and network, while \emph{control groups} (cgroups) cap its CPU time and
memory. The result is a lightweight, per-submission jail that can run untrusted code, kill
it when it exceeds a limit, and report the time and memory it consumed---while guaranteeing
that the worst a malicious submission can do is hit a limit and be terminated. Its
provenance at the IOI, where it grades the submissions of the world's strongest
high-school programmers under adversarial conditions, is exactly why it is trustworthy: it
is proven technology, not an improvisation.

\subsection{Judge0 as execution-as-a-service}

\texttt{isolate} is a sandbox, not a service. Wrapping it in a language-aware,
network-accessible API is the role of Judge0~\cite{judge0,dosilovic2018judge0}, the most
prominent open-source ``code execution as a service.'' Do\v{s}ilovi\'{c} and
Me\v{s}trovi\'{c}~\cite{dosilovic2018judge0} describe Judge0 as a system that accepts a
\texttt{\{language, source\_code, stdin\}} payload over HTTP, compiles and runs the code in
\texttt{isolate}, and returns \texttt{\{stdout, stderr, status, time, memory\}}. It supports
dozens of language runtimes behind a uniform REST interface and enforces per-submission
resource limits through the underlying sandbox.

APEX delegates all code execution to Judge0. This is a deliberate architectural choice,
not a convenience: the most security-critical component in the entire system---the one that
runs untrusted code---is the one the team did \emph{not} write, precisely because writing a
sandbox is how a project gets compromised. The whole of APEX's interaction with the sandbox
is confined to a single file, \texttt{internal/judge0/client.go}, and reduces to one
synchronous HTTP call:

\begin{lstlisting}[language=Go,caption={APEX's entire interaction with the sandbox: one synchronous POST to Judge0, which runs the code in isolate.},label={lst:bg-judge0-call}]
url := BaseURL +
  "/submissions?wait=true&base64_encoded=false" +
  "&fields=stdout,stderr,compile_output,status,time,memory"
client := &http.Client{Timeout: 30 * time.Second}
resp, err := client.Post(url, "application/json", bytes.NewReader(body))
\end{lstlisting}

Several properties of this integration are worth noting because they reflect the trade-offs
the literature anticipates. The call is \emph{synchronous}: \texttt{wait=true} blocks the
request until Judge0 finishes, guarded by a 30-second client timeout, avoiding the polling
loops or webhook callbacks that a higher-throughput design would need. The base URL defaults
to the public \texttt{https://ce.judge0.com} instance but is overridable via the
\texttt{JUDGE0\_URL} environment variable; the public endpoint is rate-limited and intended
for evaluation, so a real deployment self-hosts Judge0. APEX maps six language names to
Judge0's numeric language identifiers and pre-validates the language before ever making the
call. Notably, APEX treats Judge0 as an \emph{executor} rather than a \emph{judge}: it counts
both Judge0's ``Accepted'' (status 3) and ``Wrong Answer'' (status 4) as ``the code ran,''
then re-compares the output itself with \texttt{normalizeOutput}, because Judge0's own diff
is byte-exact and would reproduce the very trailing-newline unfairness the normalisation
exists to eliminate. This separation of ``run it safely'' from ``decide if it is correct''
is the concrete embodiment of the sandbox/scaffold boundary that C\'{a}ceres et
al.~\cite{caceres2014online} identify as characteristic of durable designs.

\section{Learning-Management Systems and Quiz Tooling}\label{sec:bg-lms}

\subsection{The LMS lineage}

The other parent of APEX is the learning-management system. Moodle~\cite{moodle2002},
first released in 2002, established the template: an open-source PHP application bundling
course pages, announcements, a quiz module, and a gradebook into a single self-hostable
product. Canvas~\cite{canvas} refined the same form factor for large institutions with a
polished interface and an extensive integration ecosystem. The defining strength of the LMS
lineage is \emph{breadth}: these systems cover the entire course lifecycle, from enrolling
students through delivering content to recording final marks.

That breadth is also the source of their limitation for programming courses. LMS quiz
modules natively handle multiple-choice, matching, and short-answer questions, but
programming-style assessment is not part of the core model. Moodle's quiz engine, for
instance, supports MCQ and numeric answers directly, but grading code requires a plugin.
The consequence, documented across the assessment surveys~\cite{douce2005autoassess,%
ihantola2010autograding}, is a workflow in which coding, multiple-choice, and written
assessment live in separate tools with separate gradebooks, and the instructor performs the
integration by hand.

\subsection{Coding on top of the LMS}

The dominant answer to ``how do I grade code inside my LMS'' is
CodeRunner~\cite{coderunner}, a mature Moodle question-type plugin that runs student code
against instructor-defined tests---typically in the unit-test style---inside a sandbox, and
reports a per-test result within the Moodle quiz. CodeRunner is capable and widely deployed,
and it is the closest LMS-native analogue to APEX's coding questions. It inherits both the
strengths of Moodle (rich rostering, gradebook integration, an established administrative
model) and its costs (a full Moodle installation, PHP hosting, plugin and sandbox
configuration, and a coding question that remains a bolt-on to a quiz engine designed around
static question types).

At the other pole sits HackerRank~\cite{hackerrank}, a commercial, cloud-hosted platform
built judge-first. It offers a polished in-browser coding experience, a large problem bank,
and strong auto-grading, and it is used both for competitive practice and for technical
hiring. Its trade-offs are the mirror image of the LMS lineage: mixed-format exams that
combine coding with substantial written and multiple-choice sections in a single sitting are
not its centre of gravity, it is neither self-hostable nor free, and its rostering and
course model are lighter than an LMS provides.

APEX sits deliberately between these poles. Rather than bolting code execution onto a
quiz engine or bolting a course model onto a judge, it treats the three answer formats
(coding, MCQ, written) as first-class citizens of a \emph{single} exam attempt: one timer,
one autosave path, one submission endpoint, and one grading pipeline that dispatches to
\texttt{Grade}, \texttt{GradeMCQ}, or \texttt{GradeWritten} by problem type before
re-aggregating a single weighted attempt score. This unification, developed in
\Cref{ch:ch07}, is the feature that neither the LMS-plus-plugin nor the judge-first
lineage offers cleanly.

\subsection{Manual grading and mixed formats}

A working classroom needs to grade prose, not only code. Written answers cannot be graded
by a diff, so any honest mixed-format system needs a manual-grading path. LMS platforms
have mature manual-grading UIs; pure judges generally do not. APEX carries a manual path as
a first-class part of the pipeline: written answers, and multiple-choice questions for which
no correct option was configured, are marked \texttt{pending\_review} and enter a queue that
teachers drain via \texttt{GET /teacher/grading/pending} and \texttt{PUT /submissions/:id/%
grade}. Crucially, \texttt{pending\_review} submissions are excluded from both the numerator
and denominator of the attempt score until a human grades them, so an ungraded essay never
makes a student's visible score look like a failure before the teacher has looked at it.
This blending of automatic and manual grading behind one attempt score is precisely the
integration the assessment surveys found missing between the two lineages.

\section{Code-Similarity and Plagiarism Detection}\label{sec:bg-plagiarism}

Adjacent to grading, and important to its fairness, is the detection of copied
submissions. The two canonical open systems are MOSS (Measure Of Software Similarity)~%
\cite{moss} from Stanford and JPlag~\cite{prechelt2002jplag}. Both ingest a corpus of
student source files and report clusters of suspiciously similar submissions, but they do
so structurally rather than textually: they tokenise each program and compare token
sequences, so that superficial obfuscations---renaming variables, reordering independent
statements, reformatting whitespace, inserting dead comments---do not defeat detection.
MOSS uses a document-fingerprinting technique (winnowing) to compare token streams at scale,
while JPlag performs a greedy tiling match over token strings; Prechelt et
al.~\cite{prechelt2002jplag} report that JPlag reliably identifies plagiarised pairs even
under such transformations.

Two properties of this work matter for APEX's scope. First, similarity detection is
\emph{corpus-relative}: it answers ``which of these submissions resemble each other,'' not
``is this one submission correct,'' and so it is orthogonal to the per-submission grading
pipeline. Second, it is typically an \emph{offline, teacher-triggered batch} action run over
a completed cohort rather than an inline check at submission time. APEX does not currently
ship an integrated similarity check; it is identified as future work in \Cref{ch:ch13}. The
architectural boundary APEX already draws---storing each submission's source, language, and
per-test results as durable rows---means a MOSS- or JPlag-style batch export is a natural
future addition rather than a redesign. Proctoring (webcam capture, screen locking) is a
separate and ethically contested body of work that APEX deliberately does not pursue; at
most a future iteration would surface non-invasive signals such as tab-blur or paste events
to the teacher rather than enforcing them.

\section{Gap Analysis and Positioning}\label{sec:bg-gap}

The surveys and systems above converge on a consistent picture. The assessment literature
identifies, and has identified for two decades, a persistent gap between \emph{course}
tooling and \emph{grading} tooling~\cite{douce2005autoassess,ihantola2010autograding}. The
LMS lineage owns the course but treats code as a bolt-on; the judge lineage owns code
execution but travels light on rostering, mixed formats, and manual grading. Bridging the
two has been the subject of a long tail of institutional projects, but the two dominant
production answers---an LMS with a coding plugin, or a commercial cloud judge---each inherit
the limitations of their lineage. Four specific gaps motivate APEX's design.

\begin{enumerate}[leftmargin=1.4em]
  \item \textbf{Mixed formats in one sitting.} No mainstream option treats coding, MCQ,
        and written answers as equal first-class citizens of a single timed attempt with a
        single score. LMS quizzes bolt code on; judges under-serve prose. APEX makes the
        unified attempt the centre of the design.
  \item \textbf{One grading pipeline.} Where mixed grading exists, it is usually the sum of
        independent subsystems. APEX runs one dispatch-and-aggregate pipeline behind all
        three formats, with a single weighted attempt score and a single ``Exam Graded''
        notification.
  \item \textbf{Self-hostability without operational weight.} Moodle is self-hostable but
        heavy; Canvas and HackerRank are effectively cloud services. APEX ships as one Go
        binary plus PostgreSQL plus an external Judge0, deployable via a four-service Docker
        Compose file, with the security-critical sandbox delegated to a proven external
        component.
  \item \textbf{Cost.} The self-hostable, no-licence-fee option (Moodle plus CodeRunner)
        is operationally demanding; the low-friction options are commercial. APEX targets
        the empty quadrant: free and self-hostable and low-friction.
\end{enumerate}

\Cref{tab:bg-comparison} summarises how APEX compares against four representative
alternatives across the dimensions that a single-instructor or single-department deployment
actually weighs. The assessment is qualitative and, in the ``mixed formats'' and
``simplicity'' rows, necessarily subjective; it is meant to make the positioning legible,
not to claim measurement. The honest caveats belong in the same breath: the four
comparators are mature, widely deployed, and battle-tested at institutional scale, whereas
APEX is a graduation project with, by its own account, no backend unit tests, no rate
limiting on authentication, and a fire-and-forget grading path with no durable work queue
(all detailed in \Cref{ch:ch08,ch:ch09}). The comparison claims a favourable \emph{shape},
not superiority on every axis.

\begin{table}[H]
  \centering
  \caption{APEX against representative alternatives. ``Yes/No'' is a capability judgement;
  cost reflects licensing and hosting for a small deployment. Assessments in the mixed-format
  and simplicity rows are subjective summaries of the discussion above.}
  \label{tab:bg-comparison}
  \small
  \renewcommand{\arraystretch}{1.3}
  \begin{tabularx}{\linewidth}{L{3.5cm} C{1.9cm} C{2.1cm} C{1.7cm} C{1.9cm} X}
    \toprule
    \textbf{Dimension} & \textbf{APEX} & \textbf{Moodle + CodeRunner} & \textbf{Canvas} &
    \textbf{Hacker\-Rank} & \textbf{Raw Judge0} \\
    \midrule
    Self-hostable
      & Yes & Yes & No (cloud) & No (cloud) & Yes \\
    Mixed coding + MCQ + written in one sitting
      & Yes, first-class & Partial (code is a plugin) & Partial (no code) & Weak & No \\
    Unified grading pipeline
      & Yes, one path & Per-subsystem & Per-subsystem & Code only & Code only \\
    Live-coding editor (Monaco, in-browser)
      & Yes & Basic & Basic & Yes & None (API only) \\
    Manual grading queue for prose
      & Yes & Yes & Yes & Weak & None \\
    Class rosters / invite codes
      & Yes & Yes & Yes & Partial & None \\
    Sandbox for untrusted code
      & Delegated to Judge0/\texttt{isolate} & Own sandbox config & External tool & Built-in &
        \texttt{isolate} \\
    Deployment simplicity
      & One Go binary + Postgres + Judge0 & Full PHP + plugin + sandbox & N/A (hosted) &
        N/A (hosted) & Sandbox only, no scaffold \\
    Cost
      & Free, self-hosted & Free, heavy ops & Paid licence & Paid & Free (self-host) \\
    \bottomrule
  \end{tabularx}
\end{table}

The pattern the table makes visible is the one the surveys predict. The LMS-derived
options (Moodle, Canvas) are strong on classroom workflow and manual grading but treat code
as a secondary format; the judge-derived options (HackerRank, raw Judge0) are strong on code
execution but thin on rostering, mixed formats, and prose grading. APEX's contribution is
not a new sandbox or a new grading algorithm---it reuses \texttt{isolate} through Judge0 and
uses textbook I/O diffing---but the \emph{integration}: unifying three answer formats and a
full classroom workflow behind one grading pipeline in a single self-hostable artifact,
which is the empty region between the two established lineages. The chapters that follow
document how that integration is engineered: the requirements it entails
(\Cref{ch:ch03}), the architecture that realises it (\Cref{ch:ch04}), and the grading
engine at its core (\Cref{ch:ch07}).
